\section{Ultrapowers and preliminary inputs}\label{n:sec:framework}

We begin a separate nonunital construction, using the hypotheses of Theorem~\ref{n:thm:main}. The immediate goal is a unital tracial quotient $D=A^\omega$ with a specified family of limit traces. Compactness of the normalized trace base supplies its unit; the later projection and model constructions take place in this quotient.

Throughout the proof, $A$ satisfies the hypotheses of
Theorem~\ref{n:thm:main}. Set $T:=T_1(A)$. In particular, every nonzero densely finite trace is
a positive scalar multiple of a unique member of $T$. Simplicity makes
all members of $T$ faithful. Projectionlessness implies that $A$ is
non-elementary and has no nonzero finite-dimensional quotient.

Fix a free ultrafilter $\omega$ on $\N$. For a trace $\sigma$, write
$\|x\|_{2,\sigma}=\sigma(x^*x)^{1/2}$, and use a set of traces in the
subscript to denote the supremum. Set
\[
 D=\Aw=\ell^\infty(A)/I_T,\qquad
 I_T=\{(a_n):\lim_\omega\|a_n\|_{2,T}=0\}.
\]
The tracial inequalities make $I_T$ a closed ideal. Let $\Lambda$ be the family
of limit traces on $D$, namely the traces given by
\[
 \tau([(a_n)])=\lim_\omega\sigma_n(a_n),\qquad \sigma_n\in T.
\]
We regard $A$ as a subalgebra of $D$ through constant sequences.

\begin{lemma}\label{n:lem:unit}
The algebra $D$ is unital, $\Lambda$ consists of tracial states, and every
$\tau\in \Lambda$ restricts to an element of $T$. Moreover,
\[
 \|[(a_n)]\|_{2,\Lambda}=\lim_\omega\|a_n\|_{2,T},
\]
and this norm detects zero in $D$.
\end{lemma}
\begin{proof}
Choose an increasing sequential contractive approximate unit $(u_n)$.
The continuous functions $\sigma\mapsto\sigma(u_n)$ increase to $1$ on
compact $T$, so Dini's theorem gives uniform convergence. In a tracial
von Neumann representation, for any bounded sequence $(a_n)$,
\[
 \|(1-u_n)a_n\|_{2,\sigma}^2
 \le\|a_n\|^2\sigma((1-u_n)^2)
 \le\|a_n\|^2(1-\sigma(u_n)).
\]
The same estimate holds on the right. Thus $(u_n)$ represents $1_D$.
Its value under every limit trace is $1$. Uniform convergence on the
constant approximate units also shows that the restriction of each limit
trace to $A$ has norm one. The norm identity follows by choosing a
trace at each coordinate within an arbitrarily small error of the
supremum. Its last assertion is the definition of the quotient.
\end{proof}

We will repeatedly use the following consequence. Bounds on finitely many
continuous expressions in the limit-trace values of bounded tuples give
uniform coordinate bounds along $\omega$, with any additional positive
slack. If such a bound failed, choosing witnessing traces on the offending
coordinates would give a contradicting limit trace. Positive contractions
and self-adjoint contractions in a $C^*$-quotient have lifts of the same
kind, by continuous functional calculus.

We record the external inputs used before the final absorption step.
A completely positive map is \emph{order zero} if it preserves
orthogonality of positive elements. All order-zero maps below are
contractive. Matrix order-zero maps lift across
arbitrary $C^*$-quotients \cite[Proposition 1.8]{CETW}. For such a map
$\theta$ from a unital algebra there is a supporting homomorphism $\pi$
in the bidual of its target such that
\[
 \theta(d)=h\pi(d),\qquad h=\theta(1),\qquad [h,\pi(d)]=0.
\]
If $f\in C([0,1])$ and $f(0)=0$, then $f(h)\pi(d)$ is in the target
algebra \cite[Theorem 3.3 and Corollary 4.2]{WZ}.

A homomorphism from a finite direct sum of matrix algebras into a quotient
consequently has an order-zero lift. Indeed, for two orthogonal image units
$e_1,e_2$ in the quotient, lift
$e_1-e_2$ as a self-adjoint contraction. Its positive and negative parts
generate orthogonal hereditary subalgebras mapping onto the corners at
$e_1,e_2$. Recurse through the summands and lift each matrix map in its
hereditary target. The ranges are orthogonal, so their sum is contractive
and order zero.

For $\sigma\in T$, let $M_\sigma=\pi_\sigma(A)''$. Kaplansky density gives
a surjection
\begin{equation}\label{n:eq:kaplansky}
 \ell^\infty(A)\longrightarrow M_\sigma^\omega.
\end{equation}
To see surjectivity, approximate the $n$th entry of a bounded sequence in
$M_\sigma$ by an element of the same ball of $A$, with $2$-norm error less
than $1/n$. This works also for the unit of $M_\sigma$. Consequently,
the inclusion of any finite-dimensional unital subalgebra of $M_\sigma$
into $M_\sigma^\omega$ has a coordinatewise order-zero lift to $A$.
This same quotient and lifting argument appears in
\cite[Proof of Lemma 4.4]{CETW}. Finally, nuclearity and Connes' injectivity theorem
make $M_\sigma$ hyperfinite, with separable predual, for every $\sigma$,
including nonextreme traces \cite{Connes}; see also
\cite[Proof of Lemma 4.5]{CETW}.

For a positive element $b$ in an ambient $C^*$-algebra $C$, write
$\Her(b)=\overline{bCb}$.
All trace calculations involving supports or supporting homomorphisms
may be made in the relevant tracial von Neumann representation. Matrix
extensions always use $\Tr$, not the normalized matrix trace.

